\documentclass[10pt]{article}
% Official TMLR style files (JmlrOrg/tmlr-style-file); see STYLE_PROVENANCE.json.
% Anonymous submission: keep \usepackage{tmlr}. Use [preprint] for arXiv, [accepted] for camera-ready.
\usepackage{tmlr}
\usepackage{amsmath,amssymb}
\usepackage{booktabs}
\usepackage{graphicx}
\usepackage{hyperref}
\usepackage{url}
\usepackage{placeins}
\graphicspath{{generated/}{../figures/}}

\newcommand{\R}{\mathbb{R}}
\newcommand{\ip}[2]{\langle #1,#2\rangle}

\title{What Does an Artist Name Add? Separating Shared and\\Painter-Specific Responses in Text-to-Image Generation}

\author{\name Anonymous authors \email anonymous@example.com \\
      \addr Paper under double-blind review}

\def\month{MM}
\def\year{YYYY}
\def\openreview{\url{https://openreview.net/forum?id=XXXX}}

\begin{document}

\maketitle

\begin{abstract}
Artist-style prompting is often evaluated by how much closer images generated with
an artist's name move to that artist's works. This proximity gain combines two
responses: a change shared by every artist name in the prompt set, and differences
between names that may or may not match the differences between the artists. We
separate them with a controlled design in which fixed scene descriptions are
rendered without a painting instruction, with a generic oil-painting instruction,
and in the style of Monet, Sisley, Pissarro or C\'ezanne, and with estimators that
use independent repeated generations to remove sampling noise from squared
magnitudes. In 1,008 images from six commercial text-to-image configurations,
66.7--88.4\% of the change that a painter name adds beyond the generic instruction
is shared by all four names in 31 interpretable image features. Name-specific
differences are present and aligned with the reference painter differences in every
configuration, but the shared change is 1.4--5.3 times as large as the entire
spread of the four painters' reference means. In CLIP and CSD embeddings of the
same images, the shared change supplies 73.2--83.8\% and 54.2--79.7\% of the gain
in similarity to the prompted painter's prototype. Proximity, agreement with the
reference painter differences, and recognition of the prompted name therefore rank
the configurations differently. A separate collection of 2,000 SD-Turbo images
shows the same majority-shared pattern overall but not in texture features. We
recommend that artist-style evaluations include a generic-style control and report
this decomposition next to any proximity score.
\end{abstract}

\section{Introduction}\label{sec:intro}

Artist names are among the most widely used style controls in text-to-image
generation, and their effect is measured in several ways. Style descriptors score
generated images by their similarity to an artist's works
\citep{somepalli2024style}; recognition benchmarks ask whether the prompted artist
can be identified from the image alone \citep{su2025prompted,moayeri2025rethinking};
and studies of style mimicry, protection and removal ask whether an artist's style
is reproduced, blocked or erased
\citep{shan2023glaze,honig2025adversarial,gandikota2023erasing,zhang2024unlearncanvas}.
A common readout across these settings is proximity: images generated ``in the
style of Monet'' are closer to Monet's paintings than images generated without the
name.

Proximity gain is ambiguous. When a prompt set contains related painters, a name
can move every output toward the group as a whole, for example toward a generic
Impressionist appearance, without separating the painters from each other. Such a
shared movement raises similarity to each painter's works, including the prompted
painter's. Whether a generator reproduces what distinguishes one painter from
another is a different question, and it requires a different measurement.

We make the distinction measurable. Our design fixes the scene text and varies
only the style clause: no painting instruction, a generic oil-painting
instruction, or an oil-painting instruction in the style of one of four related
painters, Claude Monet, Alfred Sisley, Camille Pissarro and Paul C\'ezanne.
Averaging the four named conditions splits each name's change into a component
shared by all four names and a between-name component, the departures from that
average. Only the between-name component can reproduce differences among painters.
Two independent generations per condition allow squared magnitudes to be
estimated without the positive bias of sampling noise, by multiplying independent
repeats instead of squaring one noisy estimate. We then compare the between-name
component with the differences among the four painters' digital reference
collections, and trace both components through proximity and recognition in
learned embeddings.

The main collection contains 1,008 images: six text-to-image configurations
(GPT Image 1, GPT Image 2, two GPT Image 2.5 variants named Flare and Sunburst,
Nano Banana 2 and FLUX.2 Max), 14 scenes, six clauses and two repeats. We measure
31 interpretable color, spatial and texture features and, as a representation
check, CLIP \citep{radford2021} and CSD \citep{somepalli2024style} embeddings of
the same images. A separate collection of 2,000 SD-Turbo images from an earlier
study, with a different prompt template and matched seeds, provides a check
outside this design. Our findings are:
\begin{itemize}
\item \textbf{Most of what a painter name adds is shared.} Beyond the generic
oil-painting instruction, 66.7--88.4\% of the squared change added by a name is
common to all four names in the 31 features. Name-specific differences exist, and
in every configuration they point along the reference painter differences, but the
shared change is 1.4--5.3 times the entire spread of the four reference painter
means, whereas the name-specific differences are 0.6--2.1 times that spread
(Section~\ref{sec:shared}).
\item \textbf{The same holds in learned embeddings.} The shared change supplies
73.2--83.8\% (CLIP) and 54.2--79.7\% (CSD) of the gain in similarity to the
prompted painter's reference prototype (Section~\ref{sec:learned}).
\item \textbf{Different readouts rank the configurations differently.} The
configuration with the largest proximity gain is not the one whose prompted names
are most recognizable, and the configuration with the lowest error against the
reference painter differences in the 31 features has the least recognizable names
(Section~\ref{sec:readouts}).
\item \textbf{The pattern is not universal across feature families.} In the
SD-Turbo collection the shared change is again a majority overall, 64.2\%, but not
in texture features, 36.6\% (Section~\ref{sec:sdturbo}).
\end{itemize}

The algebra behind these results is elementary; the contribution is the controlled
measurement and what it shows about evaluation. We measure digital image
properties and embeddings of reproductions, not perceived painter resemblance, and
our four painters form one related group. A large shared component is therefore
not by itself a failure: part of it may be appropriate family resemblance. The
point is that proximity gain cannot tell the two components apart, while the
decomposition can. Section~\ref{sec:limitations} states the limits of the evidence.

\section{Related work}\label{sec:related}

\paragraph{Artist style similarity and recognition.}
\citet{somepalli2024style} train contrastive style descriptors (CSD) and score
generated images against averaged artist prototypes, including prompts that
constrain content. \citet{su2025prompted} build a benchmark for identifying the
prompted artist from generated images across 110 artists, with content-controlled
name substitutions, and compare feature similarity, style descriptors and trained
classifiers. \citet{moayeri2025rethinking} decide whether an artist's style is at
risk of copying by testing whether the artist's own works are recognizable, and
find that about a fifth of 372 prolific artists qualify. \citet{frochte2026csd}
shows that raw CSD cosine fails to separate same-artist from different-artist pairs
for a substantial fraction of a 91-artist corpus, proposes a corrected readout, and
also evaluates top-1 and top-5 recognition of prompted generations. These works
establish that prompted artists are often recognizable and that raw style
similarity requires validation. We add a generic-style control and measure how much
of an artist-name response is shared across names, and how that shared part enters
proximity, agreement and recognition readouts.

\paragraph{Style mimicry, protection and removal.}
Glaze \citep{shan2023glaze} perturbs artworks to disrupt fine-tuned mimicry, and
\citet{honig2025adversarial} show with a user study that such protections can be
bypassed. Concept erasure \citep{gandikota2023erasing} and the UnlearnCanvas
benchmark \citep{zhang2024unlearncanvas} evaluate removing artistic styles from
diffusion models, and \citet{verma2025imitation} estimate how many training images
a model needs before it imitates an artist. Their readouts include classifier
recognition, embedding similarity and human judgment of whether a style is
present. Our results concern the prompting of pretrained services rather than
fine-tuning or erasure, but they suggest a question to ask of any such readout:
how much of the measured effect is a change that every artist name would produce?

\paragraph{Expert and quantitative studies of art.}
\citet{deliege2025} assess stylistic imitation with expert raters, and
AI-Pastiche \citep{asperti2025pastiche} and AI-WikiArt \citep{fu2025aiwikiart}
study the authenticity and attribution of generated paintings. Quantitative art
research has long described paintings by color, composition and image statistics
\citep{kim2014,sigaki2018,lee2020}, and recent work traces historical change in
learned embeddings \citep{kim2026}. Our features follow this tradition; we use
them to compare response components, not to define style.

\paragraph{Evaluation methodology.}
Generative evaluation separates fidelity from diversity \citep{naeem2020}, and
conditional metrics distinguish variation within a condition from relations among
condition means \citep{benny2021conditional}. Our squared-magnitude estimator is
the cross-validated distance used in representational similarity analysis
\citep{diedrichsen2017representational}. Image processing can change evaluation
features substantially \citep{parmar2022}, which motivates our source-quality
checks.

\section{Design and measurements}\label{sec:design}

\paragraph{Painters and reference collections.}
The reference collection contains 649 digital reproductions: 297 by Monet, 106 by
Sisley, 141 by Pissarro and 105 by C\'ezanne, with Wikimedia Commons source records
and Wikidata work identities \citep{commonsglam,vrandecic2014wikidata}. A separate
221-work development panel of the same four painters sets feature scaling and is
never compared with generated images. Both panels were assembled before this
experiment. The painters are a deliberately hard case: three Impressionists who
often painted similar motifs, and C\'ezanne, whose work departs more clearly from
theirs. Painter means receive equal weight throughout, so the larger Monet
collection does not dominate. The panels are finite collections of photographs of
paintings; lighting, calibration and conservation state enter the measurements.

\paragraph{Prompts.}
Fourteen fixed outdoor scene descriptions cover water, built, land and mixed
compositions (Appendix Table~\ref{tab:scenes}). Each is crossed with six clauses:
none; ``Render as an oil painting.''; and ``Render as an oil painting in the style
of [painter].'' for each painter's full name. The clause precedes the scene
description and every prompt ends with ``No text or frame.'' Nothing else varies.

\paragraph{Generators.}
We requested six configurations through one image-generation gateway: GPT Image 1,
GPT Image 2 and two GPT Image 2.5 variants (Flare and Sunburst) at medium quality,
Nano Banana 2 at 1K resolution, and FLUX.2 Max. All return
$1024\times1024$ images. Each model--scene--clause cell received two separately
requested outputs, 1,008 images in total, collected on 10 September 2026 (UTC) in a
randomized order fixed before collection. The service responses do not echo model
identifiers, so the labels identify requested configurations rather than verified
checkpoints (Appendix~\ref{app:provenance}). Figure~\ref{fig:examples} shows one
scene under all six clauses.

\begin{figure}[t]
\centering
\includegraphics[width=\linewidth]{specificity_original_generated.pdf}
\caption{One scene description (a river bending past a gravel bank, three trees on
the far bank and a low hill under an overcast sky) under all six clauses in each
configuration (first repeat, selected by a fixed rule before inspection). The top
row shows one reference painting per painter for orientation; references are not
matched to the scene. Compare the four named outputs with the two controls on the
right, then compare the names with each other.}
\label{fig:examples}
\end{figure}

\paragraph{Measurements.}
The primary representation has 31 coordinates chosen to change independently: 11
color features (lightness, chroma, hue concentration and entropy, and color
differences at three spatial scales), 8 spatial features (spectral slope and
anisotropy, gradient statistics, orientation entropy and balance, and quadrant
variation) and 12 texture features (wavelet energies, local-binary-pattern
entropies and local contrast). Appendix~\ref{app:features} defines them. Images
are converted to sRGB and resized to a 512-pixel short side. Each coordinate is
centered and scaled by the development panel's painter-weighted median and
interquartile range, so one unit is one development-panel IQR. For a
representation check we also embed every generated and reference image with CLIP
ViT-L/14 \citep{radford2021} and the released CSD ViT-L checkpoint
\citep{somepalli2024style}, using each encoder's native $224\times224$ center crop
and unit-normalized 768-dimensional outputs (Appendix~\ref{app:learned}).

\section{Separating shared and painter-specific responses}\label{sec:method}

\subsection{Shared and between-name change}

Let $z_{sak}\in\R^{31}$ be the standardized features of one configuration's image
for scene $s=1,\dots,S$ ($S=14$), clause $a$ and repeat $k\in\{1,2\}$. The clause
is $a=\mathrm{f}$ (artist-free), $a=\mathrm{g}$ (generic oil painting) or one of the
four painters $a=1,\dots,4$. Average over scenes within each repeat,
$\bar z_{ak}=S^{-1}\sum_s z_{sak}$, and write the change from a baseline
$b\in\{\mathrm{f},\mathrm{g}\}$ to painter $a$ as $h_{ak}=\bar z_{ak}-\bar z_{bk}$.
The four changes split into their mean and the departures from it:
\begin{equation}
h_{ak}=c_k+e_{ak},\qquad c_k=\tfrac14\sum_{a=1}^4 h_{ak},\qquad
e_{ak}=\bar z_{ak}-\tfrac14\sum_{a'=1}^4\bar z_{a'k}.
\label{eq:split}
\end{equation}
The shared change $c_k$ moves all four named conditions together; the between-name
changes $e_{ak}$ are the only part that can differ between painters, and they do
not depend on the baseline. We measure their squared sizes as
\begin{equation}
\text{shared}=4\ip{c_1}{c_2},\qquad
B=\sum_{a=1}^4\ip{e_{a1}}{e_{a2}},
\label{eq:components}
\end{equation}
where $\ip{u}{v}$ is the inner product. If each repeat equals a fixed mean plus
noise that has mean zero and is independent between repeats, the cross-repeat
product has the squared mean as its expectation, whereas squaring either repeat
would add the noise variance. Finite-sample estimates can therefore be negative;
we retain them. We report the shared share of the total, and we write $C$ for the
shared component against the artist-free baseline and $N$ for the shared component
against the generic baseline. Since the named-minus-free change is the
generic-minus-free shift $g_k$ plus the named-minus-generic change,
$C=G+N+I$ exactly, with $G=4\ip{g_1}{g_2}$ and a signed cross term
$I$ (Appendix~\ref{app:estimators}). $N$ isolates what the painter names add beyond
the painting instruction.

To give these magnitudes a scale, let $\mu_a$ be painter $a$'s mean reference
feature vector, $r_a=\mu_a-\frac14\sum_{a'}\mu_{a'}$ its contrast with the
four-painter mean, and $H=\sum_a\|r_a\|^2$ the spread of the four painter means.
Because $B$ and $H$ are both sums over the four painters of squared departures from
a four-painter mean, $B/H$ compares the generated painter differences with the
reference painter differences in size, and $N/H$ compares the shared change with
that same yardstick.

\subsection{Agreement with the reference painter differences}

Size is not direction. To test whether the generated differences match the
reference differences, center each scene separately,
$d_{sak}=z_{sak}-\frac14\sum_{a'}z_{sa'k}$, and define
\begin{equation}
\widehat\beta=\frac1S\sum_s\frac{1}{2}\sum_{k=1}^2
\frac{\sum_a\ip{d_{sak}}{r_a}}{H},\qquad
\widehat D=\frac1S\sum_s\frac1H\sum_a\ip{d_{sa1}-r_a}{d_{sa2}-r_a}.
\label{eq:beta-d}
\end{equation}
The aligned amplitude $\beta$ projects the generated differences onto the
reference pattern: $\beta=1$ matches its size along that pattern, and $\beta>0$
indicates an overall response in the right direction without requiring agreement
for every painter pair. The error $D$ also counts differences in all other
directions. Under the repeat model above, its expectation is the normalized
squared error of the scene-conditional mean differences: $D=0$ for exact
agreement, $D=1$ when names produce no differences at all, and $D=(b-1)^2$ when the
reference pattern is reproduced at $b$ times its size. With
$\widehat Q=(SH)^{-1}\sum_{s,a}\ip{d_{sa1}}{d_{sa2}}$, the relative squared size of
the generated differences, $\widehat D=1-2\widehat\beta+\widehat Q$, so stronger
responses need not lower the error. Because every scene is compared with the same
pooled reference target, $D$ also penalizes differences that vary across scenes;
$D_{\mathrm{agg}}$ scores the scene-averaged differences instead, and
$D=D_{\mathrm{agg}}+V_{\mathrm{scene}}$ separates the two
(Appendix~\ref{app:estimators}).

\subsection{Proximity and recognition in embeddings}

In an embedding, let $g_a$ and $g_{\mathrm{g}}$ be mean unit embeddings of the
images for painter $a$ and for the generic clause, $\mu_a$ the mean reference
embedding, and bars averages over the four painters. The average gain in cosine
similarity to the prompted painter's prototype decomposes exactly as
\begin{equation}
\underbrace{\frac14\sum_a g_a^\top\mu_a-g_{\mathrm{g}}^\top\bar\mu}_{\text{proximity gain}}
=\underbrace{(\bar g-g_{\mathrm{g}})^\top\bar\mu}_{\text{shared}}
+\underbrace{\frac{H\widehat\beta}{4}}_{\text{painter-specific}},
\label{eq:prototype-gain}
\end{equation}
with $H$ and $\widehat\beta$ computed in the embedding (Appendix~\ref{app:learned}).
The shared term is movement toward the average of the four prototypes; only the
second term depends on which name produced which image. Recognition is a separate
readout: each named image is assigned to the reference prototype with the largest
cosine similarity, and we report macro accuracy over the four painters.

\subsection{Inference}

Before collection, the protocol fixed an inferential family of 21 comparisons for
the 31 features: the six aligned amplitudes and the 15 pairwise differences in $D$.
We use paired-scene Student intervals with a Bonferroni adjustment across all 21,
which target approximate 95\% simultaneous coverage if scene-level errors are
independent. Intervals condition on the reference panels, the scaling and the 14
authored scenes, and they include fixed scene heterogeneity. The protocol also
declared the artist-free shared/between-name split as an explanatory diagnostic.
The generic-baseline split, the embedding analyses, held-scene recognition and the
SD-Turbo analysis were added afterwards, each with its rules fixed in a written plan
before its outcomes were computed, but on images whose other results were known. We
report these as descriptive results with scene- or block-deletion ranges, not
confidence intervals, and attach no significance claims to them.

\section{Results}\label{sec:results}

All 1,008 images yielded valid measurements.

\subsection{Most of what a painter name adds is shared}\label{sec:shared}

% Generated by paper/tmlr/build_assets.py; do not edit by hand.
% input reports/painter_specificity_review_v3/analysis.json sha256 79fa5ca28900e852c072e2a52d35f3cb400ee053935d857729b4ad4aa3bc542c
% input reports/painter_specificity_review_v1/analysis.json sha256 2fabb78e15dfcc7b027a73170f6ad0c87af0196c754b4e6f45275269c2c961e2
\begin{table}[t]
\caption{What a painter name adds in the 31 image features. Scene-averaged, repeat-corrected squared change, in units of the reference painter spread $H=5.915$. $G$: shift from the artist-free to the generic oil-painting prompt. $N$: shift shared by the four names beyond the generic prompt. $I$: cross term, so that the shared named-minus-free shift is $C=G+N+I$. $B$: between-name differences, which do not depend on the baseline. Shared shares are $C/(C+B)$ and $N/(N+B)$; brackets give the range over the 14 single-scene deletions.}
\label{tab:decomposition}
\centering\small
\begin{tabular}{@{}lrrrrrrc@{}}
\toprule
 & \multicolumn{4}{c}{Squared change $/H$} & \multicolumn{3}{c}{Shared share (\%)}\\
\cmidrule(lr){2-5}\cmidrule(l){6-8}
Configuration & $G$ & $N$ & $I$ & $B$ & vs.\ free & vs.\ generic & deletion range\\
\midrule
GPT Image 1 & 5.04 & 5.26 & $-$0.29 & 2.12 & 82.5 & 71.3 & [70.4, 72.3]\\
GPT Image 2 & 6.28 & 4.98 & 1.06 & 2.04 & 85.8 & 70.9 & [70.0, 72.1]\\
GPT Image 2.5 Flare & 4.01 & 4.99 & 2.29 & 2.03 & 84.8 & 71.1 & [69.8, 72.3]\\
GPT Image 2.5 Sunburst & 3.08 & 3.97 & 3.28 & 1.98 & 83.9 & 66.7 & [65.9, 67.7]\\
Nano Banana 2 & 1.22 & 1.36 & 1.91 & 0.61 & 88.1 & 69.1 & [66.8, 70.9]\\
FLUX.2 Max & 3.78 & 5.34 & 6.53 & 0.70 & 95.7 & 88.4 & [86.1, 90.2]\\
\bottomrule
\end{tabular}
\end{table}


Table~\ref{tab:decomposition} decomposes the scene-averaged change. Against the
artist-free baseline, 82.5--95.7\% of the squared change from the named clauses is
shared across the four names; this includes the painting instruction itself.
Against the generic oil-painting baseline, the share that isolates what the names
add is 66.7--88.4\%, and deleting any one scene keeps it within 65.9--90.2\%.
Computing both components within each scene before averaging, rather than after,
gives even larger artist-free shares, 88.6--97.2\%.

The shared change is not a small residual of the painting instruction. Beyond the
generic clause, the names add a shared change $N$ of 1.4--5.3 times the reference
painter spread $H=5.915$, comparable to the 1.2--6.3 times $H$ of the painting
instruction itself (Figure~\ref{fig:components}). The two shifts also interact: the
cross term $I$ is negative for GPT Image 1 and positive for the other five
configurations, so their squared sizes do not simply add.

\begin{figure}[t]
\centering
\includegraphics[width=\linewidth]{fig_components.pdf}
\caption{Squared size of the change that painter names add beyond a generic
oil-painting instruction, in units of the spread $H$ of the four reference painter
means. The between-name differences $B$ are on the order of the reference spread;
the change shared by all four names $N$ is larger in every configuration. Values
are scene-averaged and repeat-corrected in the 31 features
(Table~\ref{tab:decomposition}).}
\label{fig:components}
\end{figure}

The between-name differences, by contrast, are on the scale of the reference
differences: $B$ is 0.6--2.1 times $H$. A large shared share therefore does not
mean that painter differences are missing. GPT Image 2 produces between-name
differences of 2.04 times the reference spread, and 70.9\% of its added change is
still shared. What differs across configurations is mainly the size of the shared
movement relative to the between-name differences: FLUX.2 Max combines the largest
shared change, 5.34 times $H$, with between-name differences of only 0.70 times $H$,
and so has the largest shared share, 88.4\%.

\subsection{Painter-specific differences point the right way but do not match}\label{sec:agreement}

% Generated by paper/tmlr/build_assets.py; do not edit by hand.
% input reports/painter_specificity_v2/psv2-20260911/models.csv sha256 eed121f77626d8a37235b5bb19e370e3fbd7b6bae44b559d1b866e27e9cac41c
% input reports/painter_specificity_review_v1/analysis.json sha256 2fabb78e15dfcc7b027a73170f6ad0c87af0196c754b4e6f45275269c2c961e2
\begin{table}[t]
\caption{Agreement of the between-name differences with the reference painter differences in the 31 features. $\beta$: aligned amplitude (1 matches the reference; brackets are prespecified 95\% simultaneous intervals over 21 endpoints). $Q$: squared size of the generated painter differences relative to the reference ($Q=1$ matches). $D$: repeat-corrected error of the scene-wise differences ($D=0$ exact, $D=1$ no painter distinctions); $D_{\mathrm{agg}}$ scores the scene-averaged differences instead. Bold marks the lowest value in a column, not a significant difference.}
\label{tab:agreement}
\centering\small
\begin{tabular}{@{}lcrrr@{}}
\toprule
Configuration & $\beta$ [simultaneous interval] & $Q$ & $D$ & $D_{\mathrm{agg}}$\\
\midrule
GPT Image 1 & 0.938 [0.709, 1.168] & 2.449 & 1.572 & 1.239\\
GPT Image 2 & 0.999 [0.893, 1.104] & 2.223 & 1.226 & 1.044\\
GPT Image 2.5 Flare & 0.772 [0.680, 0.864] & 2.281 & 1.738 & 1.483\\
GPT Image 2.5 Sunburst & 0.824 [0.680, 0.968] & 2.289 & 1.641 & 1.337\\
Nano Banana 2 & 0.442 [0.256, 0.629] & 0.994 & 1.109 & \textbf{0.723}\\
FLUX.2 Max & 0.470 [0.239, 0.701] & 0.741 & \textbf{0.801} & 0.761\\
\bottomrule
\end{tabular}
\end{table}


Table~\ref{tab:agreement} compares the between-name differences with the reference
differences. Every configuration responds in the reference direction: all six
aligned amplitudes have simultaneous intervals above zero, with estimates of
0.442--0.999. Direction and error nevertheless diverge. GPT Image 2 matches the
reference amplitude almost exactly, $\beta=0.999$ [0.893, 1.104], but its painter
differences are more than twice the reference size, $Q=2.223$, so most of that size
lies off the reference pattern and its error is $D=1.226$. FLUX.2 Max has a smaller
aligned response, $\beta=0.470$, and smaller differences, $Q=0.741$, and the lowest
error estimate, $D=0.801$; its unadjusted 95\% interval [0.586, 1.016] includes the
no-difference value of 1, so no configuration is shown to beat a generator that
makes no painter distinctions at all. Of the 15 prespecified pairwise error
differences, only two are resolved: Flare and Sunburst have higher error than FLUX.2
Max (Appendix Table~\ref{tab:pairs}).

The aligned response is uneven across painters. In the 31 features, Monet--Sisley
alignment ranges from $-$0.249 to 0.129 in five configurations, against 0.402 for
GPT Image 2 (Appendix Figure~\ref{fig:pairs}), and C\'ezanne carries much of the
aggregate response: omitting him lowers FLUX.2 Max's $\beta$ from 0.470 to 0.096 and
GPT Image 2's from 0.999 to 0.651 (Appendix Table~\ref{tab:coverage}).

How the error is aggregated also changes the order. Scoring the scene-averaged
differences instead of each scene gives Nano Banana 2 a lower error than FLUX.2 Max,
0.723 against 0.761, because Nano Banana 2's painter differences vary more across
scenes. Shrinking each configuration's differences by a single scalar fitted on the
other 13 scenes, a change of feature vectors rather than images, gives GPT Image 2
the lowest held-out error, 0.554 against 0.714 for FLUX.2 Max
(Appendix Table~\ref{tab:calibration}). Each is a legitimate target; they answer
different questions.

\subsection{The same split in learned embeddings}\label{sec:learned}

% Generated by paper/tmlr/build_assets.py; do not edit by hand.
% input reports/painter_learned_audit_v1/analysis.json sha256 06c77de4be133afc203dda122b2816d4f866e1389999f18bc94b3d9c668f28f7
% input reports/painter_specificity_review_v3/analysis.json sha256 79fa5ca28900e852c072e2a52d35f3cb400ee053935d857729b4ad4aa3bc542c
\begin{table}[t]
\caption{The same split in CLIP and CSD embeddings of the same images. Shared part of prototype gain: the common term's share of the named-minus-generic gain in similarity to the prompted painter's reference prototype (Equation~\ref{eq:prototype-gain}). Shared share of squared change: as $N/(N+B)$ in Table~\ref{tab:decomposition}, computed in each embedding. Monet--Sisley $\beta$: aligned amplitude for that pair alone (1 matches the reference, negative values point the opposite way).}
\label{tab:learned}
\centering\small
\begin{tabular}{@{}lrrrrrrr@{}}
\toprule
 & \multicolumn{2}{c}{Shared part of gain (\%)} & \multicolumn{2}{c}{Shared share of change (\%)} & \multicolumn{3}{c}{Monet--Sisley $\beta$}\\
\cmidrule(lr){2-3}\cmidrule(lr){4-5}\cmidrule(l){6-8}
Configuration & CLIP & CSD & CLIP & CSD & 31 features & CLIP & CSD\\
\midrule
GPT Image 1 & 74.5 & 72.7 & 76.1 & 76.5 & 0.074 & 0.301 & 0.446\\
GPT Image 2 & 73.2 & 54.2 & 76.6 & 68.8 & 0.402 & 0.532 & 0.587\\
GPT Image 2.5 Flare & 78.8 & 65.0 & 81.7 & 74.6 & $-$0.011 & 0.339 & 0.451\\
GPT Image 2.5 Sunburst & 77.7 & 67.4 & 77.6 & 69.9 & $-$0.249 & 0.385 & 0.323\\
Nano Banana 2 & 83.8 & 79.7 & 82.3 & 77.4 & $-$0.044 & 0.318 & 0.400\\
FLUX.2 Max & 81.7 & 77.6 & 80.8 & 75.1 & 0.129 & 0.299 & 0.414\\
\bottomrule
\end{tabular}
\end{table}


The decomposition is not an artifact of the hand-designed features.
In CLIP and CSD embeddings of the same images
(Table~\ref{tab:learned}), the shared term of Equation~\ref{eq:prototype-gain}
supplies 73.2--83.8\% (CLIP) and 54.2--79.7\% (CSD) of the named-minus-generic
gain in similarity to the prompted painter's prototype. The squared-change shares
beyond the generic clause are similar, 76.1--82.3\% in CLIP and 68.8--77.4\% in CSD.
The shared term is positive in every case: naming moves images toward the average
of the four painters' prototypes.

The embeddings disagree with the 31 features on one painter pair. Monet--Sisley
alignment is positive in every configuration in both embeddings, 0.299--0.532 in
CLIP and 0.323--0.587 in CSD. The weak or reversed Monet--Sisley response in the 31
features is therefore a property of that representation, not evidence that the
generators cannot distinguish these painters. As a check on the reference panels,
the same prototypes classify the 221 development works with 79.8\% (CLIP) and
79.6\% (CSD) macro accuracy.

\subsection{Proximity, agreement and recognition rank configurations differently}\label{sec:readouts}

% Generated by paper/tmlr/build_assets.py; do not edit by hand.
% input reports/painter_learned_audit_v1/analysis.json sha256 06c77de4be133afc203dda122b2816d4f866e1389999f18bc94b3d9c668f28f7
% input reports/painter_specificity_v2/psv2-20260911/models.csv sha256 eed121f77626d8a37235b5bb19e370e3fbd7b6bae44b559d1b866e27e9cac41c
\begin{table}[t]
\caption{Three readouts of the same images rank the configurations differently. Proximity: named-minus-generic gain in cosine similarity to the prompted painter's reference prototype. Recognition: macro accuracy of assigning each of the 112 named images to its prompted painter by the nearest reference prototype (chance 25\%). Agreement: error $D$ of the between-name differences in each representation (lower is better; 1 means no painter distinctions). Bold marks the best value in each column. Values are not comparable across representations.}
\label{tab:readouts}
\centering\small
\begin{tabular}{@{}lrrrrrrr@{}}
\toprule
 & \multicolumn{2}{c}{Proximity gain} & \multicolumn{2}{c}{Recognition (\%)} & \multicolumn{3}{c}{Agreement error $D$}\\
\cmidrule(lr){2-3}\cmidrule(lr){4-5}\cmidrule(l){6-8}
Configuration & CLIP & CSD & CLIP & CSD & 31 features & CLIP & CSD\\
\midrule
GPT Image 1 & 0.072 & \textbf{0.217} & 60.7 & 72.3 & 1.572 & \textbf{0.847} & 0.735\\
GPT Image 2 & 0.100 & 0.166 & \textbf{68.8} & \textbf{75.9} & 1.226 & 1.061 & 0.741\\
GPT Image 2.5 Flare & 0.091 & 0.187 & 59.8 & 62.5 & 1.738 & 1.057 & 0.819\\
GPT Image 2.5 Sunburst & 0.101 & 0.214 & 58.9 & 61.6 & 1.641 & 1.136 & 0.944\\
Nano Banana 2 & \textbf{0.112} & 0.210 & 51.8 & 50.9 & 1.109 & 1.078 & 0.999\\
FLUX.2 Max & 0.083 & 0.217 & 41.1 & 39.3 & \textbf{0.801} & 1.058 & \textbf{0.714}\\
\bottomrule
\end{tabular}
\end{table}


Table~\ref{tab:readouts} places three readouts of the same images side by side.
Nano Banana 2 has the largest CLIP proximity gain, 0.112, yet its prompted names are
recognized in 51.8\% of images, against 68.8\% for GPT Image 2, whose gain is 0.100.
In CSD, GPT Image 2 has the smallest proximity gain, 0.166, and the highest
recognition, 75.9\%. FLUX.2 Max has the lowest agreement error in the 31 features
and in CSD, 0.714, but the least recognizable names, 41.1\% in CLIP and 39.3\% in
CSD; in CLIP, GPT Image 1 has the lowest error, 0.847. These reversals follow from
the decomposition. Proximity gain is dominated by the shared term, which carries no
information about which name was used. Recognition needs large, separable
between-name differences. The error $D$ rewards differences of the right size and
direction, and it favors small differences over large misdirected ones.

A shared offset can also change recognition without changing any painter
difference. Adding one fixed translation that aligns the mean of generated named
images with the mean of the reference prototypes, fitted on the other 13 scenes
without painter labels, leaves all centered differences unchanged but moves
recognition by $-$5.4 to +12.5 percentage points in CLIP and +0.9 to +26.8 in CSD,
with means of +2.7 and +10.0 (Appendix~\ref{app:transfer}). The effect cuts both
ways: for FLUX.2 Max in CSD it corrects 39 and harms 9 of 112 decisions. Prototypes
fitted to labeled generated images outperform both reference-based rules in all 12
configuration--encoder combinations, so the prompted names are more separable within
each generator's own outputs than with respect to the historical references.

\subsection{A separate collection: the pattern does not hold for texture}\label{sec:sdturbo}

% Generated by paper/tmlr/build_assets.py; do not edit by hand.
% input reports/painter_cross_cohort_v1/analysis.json sha256 9488dfafa5779b8b065cc756804e67ad64c2841a5c62c2f9b0105f8fef13973e
\begin{table}[t]
\caption{A separate collection: 2,000 SD-Turbo images (16 scenes, 25 matched-seed blocks). The baseline already requests an oil painting, so the shared change is comparable to $N$ in Table~\ref{tab:decomposition}. Components are in units of each feature family's own reference spread $H$. The deletion range covers the 25 single-block deletions; the within-scene share computes both components in each scene before averaging.}
\label{tab:sdturbo}
\centering\small
\begin{tabular}{@{}lrrrcrr@{}}
\toprule
Features & Shared $/H$ & Between-name $/H$ & Shared (\%) & deletion range & within scene (\%) & $\beta$\\
\midrule
All 31 & 2.06 & 1.15 & 64.2 & [63.9, 64.6] & 72.4 & 0.618\\
Color (11) & 1.29 & 0.96 & 57.5 & [57.3, 57.8] & 66.6 & 0.632\\
Spatial (8) & 4.99 & 0.89 & 84.9 & [84.4, 85.3] & 81.2 & 0.310\\
Texture (12) & 0.87 & 1.50 & 36.6 & [36.2, 36.9] & 48.5 & 0.805\\
\bottomrule
\end{tabular}
\end{table}


We applied the same decomposition, with rules fixed before computing it, to all
2,000 retained images of an earlier SD-Turbo collection: 16 scenes, 25 repeat
blocks and five conditions, with the baseline prompt ``An oil painting on canvas of
\ldots'' and the names inserted as ``by [painter]''. Within a scene and block, all
five conditions share a seed, so we form products only between different blocks,
which allows arbitrary dependence among conditions within a block
(Appendix~\ref{app:sdturbo}). Because the baseline already requests an oil painting,
the shared share is comparable to $N/(N+B)$.

Across all 31 features the shared change is again the majority, 64.2\% (63.9--64.6\%
across block deletions), 2.06 times the reference spread against 1.15 for the
between-name differences (Table~\ref{tab:sdturbo}). It is a majority in color
(57.5\%) and spatial features (84.9\%) but not in texture, 36.6\% (36.2--36.9\%);
computed within scenes, the shares are 72.4\% overall and 48.5\% for texture.
SD-Turbo's aligned amplitude is 0.618, its scene-averaged error 0.917 and its
scene-wise error 1.637, and its Monet--Sisley alignment in the 31 features is
0.471. The majority-shared pattern therefore extends to a different model, template
and seed design, but it is not a property of every feature family. This collection
reuses the same references and was analyzed before in other ways; it is a
retrospective check, not an independent replication.

\subsection{Sensitivity of the reference comparison}\label{sec:sensitivity}

The shared/between-name split uses only generated images and the fixed feature
scaling, so choices about the reference collection cannot change it. The comparisons
with the reference differences can change, and we checked them against several
alternatives (Appendix Table~\ref{tab:sensitivity}). AI assistants audited all 870
reference and development reproductions and cropped peripheral artifacts such as
frames and calibration strips from 131 of them, without human verification. With
the cropped regions, FLUX.2 Max's error becomes 0.832, and 0.872 after also refitting
the feature scaling; all six aligned amplitudes remain above zero; the Sunburst--FLUX
difference is no longer resolved, while the GPT Image 1--FLUX difference becomes
resolved. FLUX.2 Max keeps the lowest uncalibrated error point estimate under every
single-scene deletion, square measurement windows, content-matched reference
targets and two alternative feature weightings.

\section{Discussion}\label{sec:discussion}

\paragraph{What the shared change is.}
The shared change moves images toward the four painters' common appearance: in both
embeddings it increases similarity to the average reference prototype. For four
related painters, part of such a movement is appropriate, since a faithful imitation
of any one of them would also move toward the others. Our data cannot say how large
that appropriate part is, because it would require the same decomposition for real
paintings against a generic reference, which does not exist here. What the data do
show is that the shared change is several times larger than the entire spread among
the four painters, and that it dominates proximity gain in every configuration and
representation we examined. Generator internals, training data and guidance could
all contribute; our intervention identifies differences between prompt conditions,
not their causes.

\paragraph{Recommendations for artist-style evaluation.}
The results suggest four practices. First, include a generic-style control and
report how much of a proximity gain is shared across artist names; the split in
Equations~\ref{eq:split}--\ref{eq:prototype-gain} requires only a set of names and
repeated generations. Second, report agreement with the reference differences
separately from proximity, including per-pair results, because a positive aggregate
response can rest on one distinctive painter. Third, state whether errors are scored
per scene or after averaging and whether contrasts are rescaled, since each choice
changed the ordering here. Fourth, treat recognition as its own readout: it rewards
separable differences and is sensitive to shared offsets that leave every painter
difference unchanged.

\section{Limitations}\label{sec:limitations}

\paragraph{Scope.} The main experiment uses four related painters, one clause
template, 14 authored scenes and two repeats per cell. The findings describe these
painters and configurations, not artists in general. A prompt set of stylistically
distant painters, or a clause naming the shared movement itself (for example
``Impressionist''), would test whether the shared share depends on painter
similarity; we have not run either.

\paragraph{Targets.} Our targets are digital reproductions measured by 31 features
and two related encoders (CSD is built on the CLIP backbone, and all four painters
appear in CSD's training tags). None of these establishes perceived painter
resemblance, and we report no human evaluation. The reference collections mix the
paintings' properties with subject choices and reproduction processes; the
AI-assisted source audit addresses visible peripheral regions only and was not
verified by humans.

\paragraph{Services and inference.} The six configurations are closed services whose
checkpoints we cannot verify. Two repeats cannot establish that repeated requests
are independent; if they share a common error component, the cross-repeat estimates
are biased, and under an assumed common correlation of 0.251 the Nano Banana 2 and
FLUX.2 Max error estimates would swap order (Appendix~\ref{app:provenance}). The
generic-baseline decomposition, embedding, recognition and SD-Turbo analyses are
retrospective analyses of images whose other results were known, and the SD-Turbo
collection shares the historical references.

\section*{Broader impact statement}
This work concerns how to interpret measurements of artist-name prompting. It may
help evaluators avoid treating proximity to an artist's collection as evidence that a
model reproduces that artist's distinguishing characteristics, which matters for
debates about style imitation. Conversely, our measurements could be misread as a
quality ranking of commercial services or as evidence about attribution, authorship
or legal questions, which they are not. We study four historical painters whose
works are in the public domain; the reference reproductions carry public-domain, CC0
or CC BY/BY-SA records, and we do not redistribute full-resolution reference images.
No living artist or human participant is involved.

\section*{Reproducibility statement}
Section~\ref{sec:design} and the appendices give the prompts, request
configurations, features, estimators and all fixed analysis rules. The supplementary
material contains the retained feature vectors, embeddings, analysis outputs and the
code that regenerates every table and figure and checks every number quoted in the
text. Generated images are not included in the supplementary archive because of its
size limit.

\bibliography{references}
\bibliographystyle{tmlr}

\clearpage
\appendix
\section{Prompts, requests and provenance}\label{app:provenance}

% Generated by paper/tmlr/build_assets.py; do not edit by hand.
% input data/manifests/painter_specificity_v2/psv2-20260911/requests.jsonl sha256 0aa60f491be08dbd48f1f018e225b22bcae2aacd724abf3371a6312d59acc679
\begin{table}[ht]
\caption{The 14 scene descriptions, with their fixed content class. Each prompt is the clause, the description and ``No text or frame.''}
\label{tab:scenes}
\centering\small
\begin{tabular}{@{}rlp{0.78\linewidth}@{}}
\toprule
\# & Class & Scene description\\
\midrule
1 & water & A broad river bends past a gravel bank. Three tall trees stand on the far bank, with a low hill behind them under a pale overcast sky.\\
2 & water & A small pond lies between reeds and a grassy slope. A wooden landing extends from the left bank, with scattered clouds reflected on the water.\\
3 & water & A narrow canal passes beneath a low stone bridge. Garden walls line one side and leafy trees line the other in soft afternoon light.\\
4 & water & A quiet inlet meets a flat sandy shore. Two small boats rest near the water and distant headlands sit beneath a light grey sky.\\
5 & built & A village street curves past modest stone houses. A low wall and a single tree occupy the foreground, with hills visible beyond the roofs.\\
6 & built & A country house stands beside an orchard. A gravel path crosses the foreground and a small shed stands on the right under a clear morning sky.\\
7 & built & A square stone tower rises behind a cluster of low tiled roofs. An open field fills the foreground, with a few shrubs along a narrow footpath.\\
8 & built & A courtyard is enclosed by two simple farm buildings. A gate opens toward distant fields and a patch of sunlight falls across the bare ground.\\
9 & land & An open meadow slopes gently toward distant wooded hills. A narrow path crosses the grass and a few irregular clouds fill the upper sky.\\
10 & land & A rocky hillside carries scattered pine trees. A broad valley extends behind the foreground rocks, with distant mountains in hazy daylight.\\
11 & land & Rows of fruit trees stretch across a gently rising field. Tall grass borders the nearest row and a dark wood lies at the horizon.\\
12 & mixed & A river passes a small village at the foot of a hill. A stone wall crosses the foreground, with open grass and two trees beside the water.\\
13 & mixed & A farmhouse stands above a shallow stream. A footbridge connects two grassy banks and an orchard spreads behind the house in diffuse daylight.\\
14 & mixed & A lakeside path passes a low stone building and a cluster of trees. Water occupies the left half and a gentle ridge crosses the background.\\
\bottomrule
\end{tabular}
\end{table}


\paragraph{Scenes and prompts.}
The 14 scene descriptions (Table~\ref{tab:scenes}) are a fixed authored panel, not a
sample of possible content. Each prompt is the clause, a space, the scene
description and ``No text or frame.'' The named clause uses the painter's full name,
for example ``Render as an oil painting in the style of Claude Monet.'' The
artist-free prompt has no clause.

\paragraph{Request configurations.}
All requests went to one image-generation gateway with a single fixed provider per
configuration and fallbacks disabled, one image per request, square aspect ratio and
no image input or seed. The four OpenAI configurations
(\texttt{gpt-image-1}, \texttt{gpt-image-2}, \texttt{gpt-image-2.5-flare} and
\texttt{gpt-image-2.5-sunburst}) requested medium quality and an opaque background;
Nano Banana 2 (\texttt{gemini-3.1-flash-image}) requested 1K resolution; FLUX.2 Max
(\texttt{flux.2-max}) requested PNG output. All 1,008 requests succeeded on the first
collection and decode to $1024\times1024$ pixels. The two repeats of a cell are
separate requests. The exact payloads and a randomized request order were fixed and
hashed before collection, which ran on 10 September 2026 between 15:34 and 18:21 UTC.
The responses do not echo model, quality or size fields, so the labels identify the
requested configurations; the underlying checkpoints and their training data cannot be
inspected. The comparison concerns these configurations, not each service's best
attainable quality or equal cost per image.

\paragraph{Repeat dependence.}
The cross-repeat estimators assume that the two repeats of a cell have independent
errors. Two repeats cannot test this. If both repeats of every cell shared an error
component with a common fraction $\rho$ of the noise power, the estimated error of
each configuration would be biased downward by an amount proportional to its
observed repeat disagreement. Recomputing all 15 pairwise differences under this
assumption, three point orderings reverse within $\rho<1$; the Nano Banana 2--FLUX.2
Max ordering reverses at $\rho=0.251$ and the GPT Image 1--GPT Image 2 ordering at
$\rho=0.246$. The actual dependence is not identified, so we treat the pairwise
orderings as conditional on independent repeats.

\FloatBarrier
\section{Feature definitions}\label{app:features}

\begin{table}[ht]
\caption{The 31 image features. Counts are scalar coordinates. Texture features
describe the digital image; they do not measure physical paint relief.}
\label{tab:features}
\centering\small
\begin{tabular}{@{}p{0.14\linewidth}p{0.42\linewidth}p{0.36\linewidth}@{}}
\toprule
Family & Measurements & Interpretation\\
\midrule
Color (11) & Lightness and chroma median and IQR; chromatic fraction; hue
concentration and entropy; color differences at three lags and their slope &
Brightness, color strength and diversity, and how quickly colors change across
the image\\[3pt]
Spatial (8) & Spectral slope, residual and anisotropy; orientation entropy;
horizontal--vertical balance; gradient median and IQR; quadrant divergence &
Coarse versus fine organization, dominant directions, edge strength and regional
variation in orientation\\[3pt]
Texture (12) & Four wavelet energies, their slope and curvature; three
local-binary-pattern entropies; three local coefficients of variation & Detail across
scales, diversity of local patterns and local luminance contrast\\
\bottomrule
\end{tabular}
\end{table}

Images are orientation-corrected and converted to sRGB using valid embedded color
profiles; unprofiled files are treated as sRGB. They are resized with Lanczos
filtering to a 512-pixel short side without upsampling, preserving aspect ratio. No
painting-region mask is applied in the primary analysis.

\paragraph{Color.} Features use CIELAB under D65 and the $2^\circ$ observer.
Lightness and chroma each contribute a median and an IQR. The chromatic fraction is
the share of pixels with $C^*\ge5$. Hue concentration is the length of the
chroma-weighted circular mean hue; hue entropy uses 24 equal bins, normalized by the
log of the bin count, with chroma weights; both are zero when fewer than 1\% of pixels
are chromatic. At lags of 1\%, 4\% and 16\% of the short side, the median CIEDE2000
difference \citep{sharma2005ciede} is taken over pooled horizontal and vertical pixel
pairs, and a log--log slope summarizes the three.

\paragraph{Spatial.} Features use linear-sRGB luminance. After a Tukey window
(parameter 0.1), 32 logarithmic radial bins cover 4--128 cycles per short side; a
Theil--Sen line gives the spectral slope and its root-mean-square residual, and
power-weighted axial concentration gives anisotropy. Scharr derivatives give the
gradient median and IQR, an 18-bin magnitude-weighted orientation entropy and the
horizontal--vertical balance. Quadrant divergence is the mean Jensen--Shannon
divergence of the four quadrants' orientation histograms from the global histogram.

\paragraph{Texture.} A four-level undecimated \texttt{db2} wavelet transform
\citep{mallat1989wavelet} gives four log mean-squared detail energies, their linear
slope and mean second difference. Uniform local-binary-pattern entropies
\citep{ojala2002lbp} use $(P,R)=(8,1),(16,2),(32,4)$ on quantized luminance. Three
median local coefficients of variation use windows of 1\%, 4\% and 16\% of the short
side (rounded up to odd widths) with a $10^{-6}$ floor on the mean.

\paragraph{Scaling.} Each coordinate is centered by the median and divided by the IQR
of the 221-work development panel (101 Monet, 36 Sisley, 48 Pissarro and 36
C\'ezanne), with each painter given equal total weight. The Euclidean metric then
weights all 31 standardized coordinates equally; Appendix Table~\ref{tab:sensitivity}
reports equal-family and covariance weightings.

\FloatBarrier
\section{Estimators}\label{app:estimators}

\paragraph{Noise correction.} Write a repeat as $x_k=\theta+\eta_k$ with fixed mean
$\theta$ and errors of mean zero that are independent between $k=1,2$. Then
$\mathbb E\ip{x_1}{x_2}=\|\theta\|^2$, whereas
$\mathbb E\|x_1\|^2=\|\theta\|^2+\mathbb E\|\eta_1\|^2$. All squared magnitudes in the
paper use such cross-repeat products. Centering across the four painters couples
the painters within a repeat but does not introduce dependence between repeats.

\paragraph{Baselines and the cross term.} Averaging scenes within repeat $k$, let
$g_k$ be the generic-minus-free shift, $n_k$ the mean named-minus-generic shift and
$c_k=g_k+n_k$ the mean named-minus-free shift. Then
\begin{equation}
\underbrace{4\ip{c_1}{c_2}}_{C}=\underbrace{4\ip{g_1}{g_2}}_{G}
+\underbrace{4\ip{n_1}{n_2}}_{N}
+\underbrace{4\bigl(\ip{g_1}{n_2}+\ip{n_1}{g_2}\bigr)}_{I}.
\end{equation}
The between-name component $B$ uses the centered named vectors and is identical for
both baselines. The reported shares are $C/(C+B)$ and $N/(N+B)$, computed when the
denominator is positive; the interaction $I$ is an algebraic term, not an identified
mechanism. The within-scene variant computes the shared and between-name components
in each scene before averaging them over scenes.

\paragraph{Error identities.} With $\theta_{sa}=\mathbb E[d_{sak}]$ under the repeat
model, $\mathbb E[\widehat D]=(SH)^{-1}\sum_{s,a}\|\theta_{sa}-r_a\|^2$. A purely shared
response has $\theta_{sa}=0$ and expected error 1, and $\theta_{sa}=b\,r_a$ gives
$(b-1)^2$. Writing $\bar\theta_a=S^{-1}\sum_s\theta_{sa}$,
\begin{equation}
D=\underbrace{\frac1H\sum_a\|\bar\theta_a-r_a\|^2}_{D_{\mathrm{agg}}}
+\underbrace{\frac{1}{SH}\sum_{s,a}\|\theta_{sa}-\bar\theta_a\|^2}_{V_{\mathrm{scene}}},
\end{equation}
so scene-wise error adds the variation of the painter differences across scenes to
the error of their average. The cross-repeat version of this identity is exact for the
retained vectors; its variation term can be negative in finite samples. Rescaling
every centered generated contrast by $c$ gives $D(c)=1-2c\beta+c^2Q$; we fit
$c=\max(0,\beta/Q)$ on 13 scenes and score the held-out scene.

\paragraph{Intervals.} For a pair of configurations, one error difference per scene
gives a paired-scene standard error; the simultaneous half-width is
$t_{13,\,1-0.05/42}$ times that standard error, and the six aligned amplitudes use
the same adjustment. The expected squared standard error includes the variance of the
fixed scene effects, so the intervals describe these 14 scenes and do not isolate the
variability of fresh requests. In 5,000 simulated experiments per scenario with the
same design, independent Gaussian, heavy-tailed heteroskedastic and scene-interaction
errors gave simultaneous coverage of 96.0\%, 97.4\% and 99.7\%, whereas a
model-specific state shared by all scenes and both repeats reduced it to 0.04\%. This
is why repeat dependence is a stated limitation.

\FloatBarrier
\section{Additional results for the 31 features}\label{app:hand}

% Generated by paper/tmlr/build_assets.py; do not edit by hand.
% input reports/painter_specificity_v2/psv2-20260911/model_contrasts.csv sha256 0d0b722750e657b83658ccffbc7cb093d7b0ae577896904c45c02ababe31665b
\begin{table}[ht]
\caption{All 15 prespecified pairwise error comparisons in the 31 features. Brackets are approximate 95\% simultaneous intervals over the 21-endpoint family; a positive difference favors configuration B. $^{*}$: interval excludes zero.}
\label{tab:pairs}
\centering\small
\begin{tabular}{@{}llc@{}}
\toprule
Configuration A & Configuration B & $D_A-D_B$ [simultaneous interval]\\
\midrule
GPT Image 1 & GPT Image 2 & 0.346 [$-$0.400, 1.093]\\
GPT Image 1 & GPT Image 2.5 Flare & $-$0.165 [$-$1.027, 0.697]\\
GPT Image 1 & GPT Image 2.5 Sunburst & $-$0.069 [$-$1.181, 1.043]\\
GPT Image 1 & Nano Banana 2 & 0.463 [$-$0.602, 1.528]\\
GPT Image 1 & FLUX.2 Max & 0.771 [$-$0.037, 1.580]\\
GPT Image 2 & GPT Image 2.5 Flare & $-$0.512 [$-$1.150, 0.126]\\
GPT Image 2 & GPT Image 2.5 Sunburst & $-$0.415 [$-$1.181, 0.350]\\
GPT Image 2 & Nano Banana 2 & 0.116 [$-$0.922, 1.155]\\
GPT Image 2 & FLUX.2 Max & 0.425 [$-$0.161, 1.011]\\
GPT Image 2.5 Flare & GPT Image 2.5 Sunburst & 0.096 [$-$0.286, 0.478]\\
GPT Image 2.5 Flare & Nano Banana 2 & 0.628 [$-$0.283, 1.540]\\
GPT Image 2.5 Flare & FLUX.2 Max & 0.937 [0.256, 1.618]$^{*}$\\
GPT Image 2.5 Sunburst & Nano Banana 2 & 0.532 [$-$0.436, 1.500]\\
GPT Image 2.5 Sunburst & FLUX.2 Max & 0.840 [0.058, 1.623]$^{*}$\\
Nano Banana 2 & FLUX.2 Max & 0.309 [$-$0.847, 1.464]\\
\bottomrule
\end{tabular}
\end{table}

% Generated by paper/tmlr/build_assets.py; do not edit by hand.
% input reports/painter_specificity_review_v1/analysis.json sha256 2fabb78e15dfcc7b027a73170f6ad0c87af0196c754b4e6f45275269c2c961e2
\begin{table}[ht]
\caption{Scene aggregation and contrast magnitude in the 31 features. $D=D_{\mathrm{agg}}+V_{\mathrm{scene}}$. Alignment is $\beta/\sqrt{Q}$. Rescaling multiplies every centered generated contrast by a nonnegative scalar fitted on the other 13 scenes and scores the held-out scene; the fitted range covers the 14 folds. Rescaling changes feature vectors, not images.}
\label{tab:calibration}
\centering\small
\begin{tabular}{@{}lrrrrcr@{}}
\toprule
Configuration & $D$ & $D_{\mathrm{agg}}$ & $V_{\mathrm{scene}}$ & $\beta/\sqrt{Q}$ & fitted scalar & $D_{\mathrm{held}}$\\
\midrule
GPT Image 1 & 1.572 & 1.239 & 0.333 & 0.600 & 0.366--0.398 & 0.645\\
GPT Image 2 & 1.226 & 1.044 & 0.182 & 0.670 & 0.438--0.459 & \textbf{0.554}\\
GPT Image 2.5 Flare & 1.738 & 1.483 & 0.255 & 0.511 & 0.329--0.349 & 0.741\\
GPT Image 2.5 Sunburst & 1.641 & 1.337 & 0.305 & 0.545 & 0.346--0.373 & 0.706\\
Nano Banana 2 & 1.109 & 0.723 & 0.387 & 0.444 & 0.401--0.533 & 0.832\\
FLUX.2 Max & 0.801 & 0.761 & 0.040 & 0.546 & 0.598--0.694 & 0.714\\
\bottomrule
\end{tabular}
\end{table}

% Generated by paper/tmlr/build_assets.py; do not edit by hand.
% input reports/painter_specificity_review_v1/analysis.json sha256 2fabb78e15dfcc7b027a73170f6ad0c87af0196c754b4e6f45275269c2c961e2
% input reports/painter_specificity_review_v3/analysis.json sha256 79fa5ca28900e852c072e2a52d35f3cb400ee053935d857729b4ad4aa3bc542c
\begin{table}[ht]
\caption{Which painter differences carry the aligned response. Components: aligned amplitude along each of the three singular components of the centered reference means (66.3\%, 20.6\% and 13.0\% of their spread). Omitting a painter recenters the other three and recomputes $\beta$. Monet--Sisley $\beta$ is shown for full frames and for central-square reference windows.}
\label{tab:coverage}
\centering\small
\begin{tabular}{@{}lrrrrrrrrr@{}}
\toprule
 & \multicolumn{3}{c}{Component} & \multicolumn{4}{c}{$\beta$ after omitting} & \multicolumn{2}{c}{Monet--Sisley $\beta$}\\
\cmidrule(lr){2-4}\cmidrule(lr){5-8}\cmidrule(l){9-10}
Configuration & 1st & 2nd & 3rd & Monet & Sisley & Pissarro & C\'ezanne & full & square\\
\midrule
GPT Image 1 & 1.248 & 0.310 & 0.356 & 1.140 & 1.072 & 0.840 & 0.389 & 0.074 & 0.038\\
GPT Image 2 & 1.156 & 0.617 & 0.800 & 1.195 & 0.972 & 0.992 & 0.651 & 0.402 & 0.408\\
GPT Image 2.5 Flare & 1.022 & 0.211 & 0.383 & 1.034 & 0.734 & 0.802 & 0.222 & $-$0.011 & 0.012\\
GPT Image 2.5 Sunburst & 1.082 & 0.250 & 0.417 & 1.116 & 0.840 & 0.762 & 0.284 & $-$0.249 & $-$0.185\\
Nano Banana 2 & 0.513 & 0.371 & 0.195 & 0.569 & 0.446 & 0.388 & 0.276 & $-$0.044 & $-$0.017\\
FLUX.2 Max & 0.666 & 0.021 & 0.183 & 0.539 & 0.511 & 0.527 & 0.096 & 0.129 & 0.132\\
\bottomrule
\end{tabular}
\end{table}

% Generated by paper/tmlr/build_assets.py; do not edit by hand.
% input reports/painter_specificity_review_v1/analysis.json sha256 2fabb78e15dfcc7b027a73170f6ad0c87af0196c754b4e6f45275269c2c961e2
% input reports/painter_reference_quality_v1/analysis.json sha256 089f636335129d6377328088ccc9fdc763488aba2a88463939322326707052b4
\begin{table}[ht]
\caption{Error $D$ under alternative targets, weightings and source corrections. Class target: title-derived water, built and land reference means for the 11 matching scenes (pooled: the same 11 scenes against the pooled target). Equal family: each feature family has equal total weight. Covariance: shrunk inverse covariance of the development panel. Cropped: AI-audited painting regions with the original scaler; refit: also refitting the development scaling. Each column has its own normalizer.}
\label{tab:sensitivity}
\centering\small
\begin{tabular}{@{}lrrrrrrr@{}}
\toprule
 & & \multicolumn{2}{c}{11 scenes} & \multicolumn{2}{c}{Weighting} & \multicolumn{2}{c}{Source correction}\\
\cmidrule(lr){3-4}\cmidrule(lr){5-6}\cmidrule(l){7-8}
Configuration & Primary & pooled & class & equal family & covariance & cropped & refit\\
\midrule
GPT Image 1 & 1.572 & 1.557 & 1.366 & 1.551 & 2.103 & 1.561 & 1.631\\
GPT Image 2 & 1.226 & 1.099 & 1.130 & 1.243 & 1.511 & 1.171 & 1.259\\
GPT Image 2.5 Flare & 1.738 & 1.726 & 1.539 & 1.765 & 1.636 & 1.680 & 1.737\\
GPT Image 2.5 Sunburst & 1.641 & 1.633 & 1.525 & 1.680 & 1.617 & 1.536 & 1.577\\
Nano Banana 2 & 1.109 & 0.987 & 0.948 & 1.203 & 1.385 & 1.076 & 1.121\\
FLUX.2 Max & 0.801 & 0.750 & 0.786 & 0.808 & 0.928 & 0.832 & 0.872\\
\bottomrule
\end{tabular}
\end{table}


\begin{figure}[ht]
\centering
\includegraphics[width=\linewidth]{specificity_artist_pairs.pdf}
\caption{Aligned amplitude and error for each of the six painter pairs (M: Monet,
S: Sisley, P: Pissarro, C: C\'ezanne), using all 31 features and normalized by the
squared reference distance of each pair. An amplitude of 1 matches the reference,
negative values point the opposite way, and an error of 1 corresponds to no
distinction. Negative errors can arise from repeat correction.}
\label{fig:pairs}
\end{figure}

\paragraph{Painter pairs and coverage.} Table~\ref{tab:coverage} shows that the
aggregate aligned response rests mostly on the first reference component, which
accounts for 66.3\% of the reference spread, and that omitting C\'ezanne removes much
of it. Stronger pair alignment does not ensure lower pair error: GPT Image 1 has weaker
Monet--Sisley alignment than GPT Image 2 but lower error for that pair
(Figure~\ref{fig:pairs}). For Flare, Sunburst and Nano Banana 2, swapping the Monet
and Sisley reference labels lowers the overall error, as it must when the pair
alignment is negative.

\paragraph{Scene aggregation and rescaling.} Table~\ref{tab:calibration} shows that
Nano Banana 2 has the largest scene variation and FLUX.2 Max the smallest, which
explains their reversal between $D$ and $D_{\mathrm{agg}}$. GPT Image 2's differences
are about twice the reference size, so shrinking them to about 0.45 of their size
removes more off-pattern difference than aligned response. Deleting any scene and
refitting the whole procedure preserves GPT Image 2's advantage after rescaling.

\paragraph{Reference sensitivity.} The content-matched target uses title-derived
water, built and land classes for the 11 non-mixed scenes. The source audit recorded
removable frames, margins, labels and calibration strips; 43 uncertain boundaries
were left unchanged. All works and features are retained under every
correction.

\FloatBarrier
\section{Learned embeddings}\label{app:learned}

\paragraph{Extraction.} We used the public checkpoints
\texttt{openai/clip-vit-large-patch14} and \texttt{tomg-group-umd/CSD-ViT-L}. CSD uses
its released style projection of the CLIP visual backbone; its release notes a
discrepancy with the published benchmark numbers, so we characterize the released
checkpoint without claiming to reproduce those numbers. Every image is
orientation- and color-corrected as above, resized bicubically to a 224-pixel short
side, center-cropped to $224\times224$ and normalized with the CLIP channel
statistics. Embeddings are unit-normalized 768-dimensional vectors; we use their
Euclidean and cosine geometry without scaling or dimension selection. Backend and
batch-size checks on fixed test cases agreed to within $1.3\times10^{-6}$ per
coordinate.

\paragraph{Prototype gain.} With unnormalized prototypes $\mu_a$ (means of unit
reference embeddings), $g_a^\top\mu_b$ is the average image-to-reference cosine
similarity. Writing $g_a=\bar g+d_a$ and $\mu_a=\bar\mu+r_a$,
\begin{equation}
\frac14\sum_a g_a^\top\mu_a-g_{\mathrm g}^\top\bar\mu
=(\bar g-g_{\mathrm g})^\top\bar\mu+\frac14\sum_a d_a^\top r_a
=(\bar g-g_{\mathrm g})^\top\bar\mu+\frac{H\widehat\beta}{4},
\end{equation}
because $\sum_a r_a=\sum_a d_a=0$. The first term cannot depend on which name produced
which images: permuting the painter labels of the generated images leaves it
unchanged and changes only $\widehat\beta$. For this reason, ranking label
permutations by prototype similarity, by $\beta$ or by $D$ gives the same order by
construction, and such permutations are not independent tests.

\paragraph{Scope.} CSD is built on the CLIP backbone and both encoders see the same
images, so their agreement is evidence about related representations, not independent
confirmation. All four painters appear in CSD's style-tag vocabulary, and the overlap
between either encoder's training data and our references is unknown. Using the
audited source regions or the development panel as the reference target changes the
shared part of the prototype gain by at most 1.6 percentage points in any
configuration.

\FloatBarrier
\section{Held-scene prompt-name recognition}\label{app:transfer}

% Generated by paper/tmlr/build_assets.py; do not edit by hand.
% input reports/painter_prototype_transfer_v1/analysis.json sha256 e1d1fbb924feb797f67e0907677511ed7c745335a4f2827e99d2f286e6be5448
\begin{table}[ht]
\caption{Held-scene recognition of the prompted painter (\%, 112 named images per configuration, chance 25\%). B: nearest reference prototype. T: the same after adding one fixed translation that aligns the generated and reference mixture means, fitted on the other 13 scenes without painter labels. G: nearest generated-image prototype fitted with painter labels on the other 13 scenes. T and G use more information than B, and G more than T.}
\label{tab:transfer}
\centering\small
\begin{tabular}{@{}lrrrrrr@{}}
\toprule
 & \multicolumn{3}{c}{CLIP} & \multicolumn{3}{c}{CSD}\\
\cmidrule(lr){2-4}\cmidrule(l){5-7}
Configuration & B & T & G & B & T & G\\
\midrule
GPT Image 1 & 60.7 & 55.4 & 80.4 & 72.3 & 73.2 & 84.8\\
GPT Image 2 & 68.8 & 66.1 & 75.0 & 75.9 & 76.8 & 94.6\\
GPT Image 2.5 Flare & 59.8 & 59.8 & 77.7 & 62.5 & 74.1 & 92.9\\
GPT Image 2.5 Sunburst & 58.9 & 62.5 & 69.6 & 61.6 & 70.5 & 90.2\\
Nano Banana 2 & 51.8 & 59.8 & 63.4 & 50.9 & 61.6 & 75.0\\
FLUX.2 Max & 41.1 & 53.6 & 71.4 & 39.3 & 66.1 & 75.0\\
\midrule
Mean change vs.\ B (pp) & & +2.7 & +16.1 & & +10.0 & +25.0\\
\bottomrule
\end{tabular}
\end{table}


For each configuration and held-out scene $s$, let $x$ be a unit embedding of a
named image from scene $s$ and $u_a=\mu_a/\|\mu_a\|$ the normalized reference
prototypes. The reference rule is $B(x)=\arg\max_a x^\top u_a$. The translation rule
first computes $t_{-s}=\frac14\sum_a\mu_a-\bar g_{-s}$, where $\bar g_{-s}$ is the
mean of the 104 named embeddings from the other 13 scenes, and predicts
$\arg\max_a(x+t_{-s})^\top u_a$. The translation changes each class score by the
constant $t_{-s}^\top u_a$ and cancels exactly under centering across painters, so
$\beta$, $Q$ and $D$ are unchanged. The generated-prototype rule uses the normalized
mean of the 26 labeled training embeddings for each painter. Held-scene repeats never
enter the fit, and the artist-free and generic images are not used. The question was
chosen after the reference-rule confusions were known; the rules were fixed before
the translated outcomes were computed. Folds overlap and scenes are authored, so these
are descriptive out-of-fold results. Across the four combinations of source view
and reference target, the mean translation gains are 2.2--4.0 points (CLIP) and
9.7--10.1 points (CSD).

\FloatBarrier
\section{The SD-Turbo collection}\label{app:sdturbo}

The collection was generated on 5 September 2026 with \texttt{stabilityai/sd-turbo}
at $512\times512$ pixels, one denoising step and guidance scale zero, for an earlier
distance analysis. It crosses 16 scenes, 25 repeat blocks and five conditions (the
baseline and four names), 400 images per condition. Within each scene and block, all
five conditions use the same seed. Its image hashes are disjoint from the main
collection.

For block vectors $v_b$, $b=1,\dots,K$ with $K=25$, the estimator
$U(v)=[K(K-1)]^{-1}\sum_{b\ne b'}\ip{v_b}{v_{b'}}$ replaces the two-repeat product.
Averaging scenes within each block, let $\delta_{ba}$ be the named-minus-baseline
shift, $c_b$ its mean over names and $d_{ba}=\delta_{ba}-c_b$. The shared component is
$4U(c)$ and the between-name component is $\sum_aU(d_a)$. Because products are formed
only between different blocks, correlation among the five seed-matched conditions
within a block is allowed. The reference means and scaling are those of the main
experiment. The within-scene variant computes the components in each scene before
averaging. Deleting each block once gives the reported ranges; they describe
sensitivity, not uncertainty.

\section{Evidence and replay}\label{app:replay}

Every analysis in this paper reads retained feature vectors or embeddings and writes a
versioned result file whose inputs are recorded by SHA-256. The tables and
Figure~\ref{fig:components} are generated from those result files, and a check script
verifies the input hashes, regenerates every table and figure byte for byte and
confirms that each number quoted in the text equals its recomputed value. These
checks establish that the reported numbers follow from the retained measurements; they
do not re-extract features from pixels, which requires the image files.


\end{document}
